\chapter{Conclusioni e sviluppi futuri}
\label{chap:conclusioni}

\section{Il problema affrontato}

% Riprendere in poche righe il problema e la domanda di ricerca dell'introduzione.

\section{Sintesi dei risultati}

% Riprendere gli obiettivi dell'introduzione e, per ciascuno, ciò che è stato
% realizzato e ciò che i risultati mostrano. Riportare soltanto conclusioni
% sostenute dal Capitolo dei risultati; non introdurre numeri nuovi.

\section{Limiti dello studio}

% Natura teorica del target PVGIS, dipendenza spazio-temporale dei campioni,
% sbilanciamento dei quattro casi input--target e assenza di etichette esterne.
% STGAN usa normalizzazione dello score e ranking sull'intero test: protocollo
% retrospettivo. Chiarire i limiti per un eventuale utilizzo online.
% Valutare separatamente l'effetto del contesto spaziale e della variazione
% del numero di parametri, senza attribuire causalità a soli confronti aggregati.
% Per ERA5 discutere selezione delle celle complete, griglia angolare, differenze
% fra soglia percentile e top-K esatto e dipendenza dei cluster dalla morfologia.
% La ricostruzione degli eventi è retrospettiva; la loro corrispondenza con
% fenomeni meteorologici reali richiede una validazione indipendente.
% La dispersione Monte Carlo dipende dal solo dropout del generatore e non è
% calibrata; il monitoraggio senza etichette orienta la selezione ma non
% sostituisce una valutazione del rilevamento. I confronti fra patch, campo
% intero e GAT non sono a parità di aggiornamenti dei pesi.
% In assenza di una verità di riferimento indipendente, l'accuratezza del
% rilevamento non è misurata.

\section{Sviluppi futuri}

% Valutazione con anomalie iniettate in posizioni note o con cataloghi esterni
% di eventi. Confronti a parità di aggiornamenti fra le varianti del detector.
% Calibrazione della dispersione dello score. Estensione ad altre aree o periodi.
% Impiego delle decisioni del detector all'interno del previsore, e non soltanto
% nell'analisi successiva.
